\documentclass[11pt,a4paper]{article}
\usepackage[margin=2.5cm]{geometry}
\usepackage{amsmath,amssymb}
\usepackage{booktabs}
\usepackage{hyperref}
\usepackage{listings}
\usepackage{xcolor}

\lstset{
  basicstyle=\ttfamily\small,
  keywordstyle=\color{blue},
  commentstyle=\color{gray},
  breaklines=true,
  frame=single,
  columns=fullflexible
}

\title{Supercell ISDF (SC\_ISDF): Algorithm and Code Implementation}
\author{GW\_double\_test / test\_SC}
\date{June 11, 2026}

\begin{document}
\maketitle

%======================================================================
\section{Introduction}
%======================================================================

In GW and COHSEX calculations, the interpolation separable density fitting (ISDF)
method approximates the exchange interaction using a small set of real-space
sampling points $\{\mathbf R_\mu\}$ with associated wavefunction values
$\zeta_\mu(\mathbf r) = \psi(\mathbf r_\mu)$.
Adaptive ISDF on the \emph{unit cell} selects these points iteratively until a
Schur-complement accuracy threshold is met.

For \emph{supercell} calculations (larger real-space FFT grids and more atoms),
re-running full adaptive ISDF is expensive.
\textbf{SC\_ISDF} reuses the unit-cell adaptive grid: it replicates the sampling
points into each translated unit-cell image inside the supercell, recomputes FFT
linear indices on the supercell grid, and then rebuilds coefficients and
$\tilde V^q$ on the supercell wavefunctions.

This document maps each algorithmic step to the MATLAB implementation in
\texttt{GW\_double\_test/service/+isdf/}.
Test cases live under \texttt{test\_profile/test\_SC/} (Si and LiH).

%======================================================================
\section{Algorithm and Corresponding Code}
%======================================================================

%----------------------------------------------------------------------
\subsection{Supercell and Unit Cell}
%----------------------------------------------------------------------

\subsubsection{Equations and Algorithm Description}

\paragraph{Real-space FFT grids.}
Let the unit cell use FFT dimensions $\mathbf N_{\mathrm{uc}} = (N_1,N_2,N_3)$
and the supercell be a $k$-fold replication
$\mathbf k = (k_1,k_2,k_3)$ in the three lattice directions:
\begin{equation}
  \mathbf N_{\mathrm{sc}} = \mathbf N_{\mathrm{uc}} \odot \mathbf k,
  \qquad
  N_{\mathrm{sc},i} = N_{\mathrm{uc},i}\, k_i .
\end{equation}
The GW code requires \texttt{FFT.get().fftgrid} on the supercell to equal
$\mathbf N_{\mathrm{sc}}$ before calling SC\_ISDF.

\paragraph{Sampling-point replication.}
Unit-cell ISDF provides $N_\mu^{\mathrm{uc}}$ points with RLU coordinates
$\mathbf R^{\mathrm{uc}}_\mu \in [0,N_{\mathrm{uc},1})\times\cdots$
(modulo periodicity).
For each image $(i_1,i_2,i_3)$ with $0\le i_j<k_j$,
\begin{equation}
  \mathbf R^{\mathrm{sc}}_{\mu,(i_1,i_2,i_3)}
  = \mathbf R^{\mathrm{uc}}_\mu
    + \bigl(i_1 N_{\mathrm{uc},1},\, i_2 N_{\mathrm{uc},2},\, i_3 N_{\mathrm{uc},3}\bigr)
  \pmod{\mathbf N_{\mathrm{sc}}}.
\end{equation}
This yields
$N_\mu^{\mathrm{sc}} = N_\mu^{\mathrm{uc}}\, k_1 k_2 k_3$ points.
Coarse and extra subsets ($N_{\mathrm{coarse}}$, $N_{\mathrm{extra}}$) are
replicated likewise.

\paragraph{Linear index on the supercell fine grid.}
Each $\mathbf R^{\mathrm{sc}}$ is mapped to a linear index on the supercell FFT mesh:
\begin{equation}
  \mathrm{lin}(\mathbf R)
  = 1 + R_1 + R_2 N_{\mathrm{sc},1} + R_3 N_{\mathrm{sc},1} N_{\mathrm{sc},2}.
\end{equation}

\paragraph{Symmetry rotation table.}
For each sampling row $a$ and crystal symmetry $S$,
\begin{equation}
  \bigl(\mathrm{R\_rot\_extra}\bigr)_{a,S}
  = \mathrm{lin}\!\bigl(\mathrm{round}(\mathbf R_a \cdot S)\bigr),
\end{equation}
using the supercell \texttt{Rgrid\_RLU} and \texttt{symmetry.get()} operations.

\paragraph{Post-replication pipeline.}
After grid replication:
\begin{enumerate}
  \item Extract $\texttt{coeff\_seper}$ on the supercell WFs at $\mathrm{lin}(\mathbf R_\mu)$
        (\texttt{gen\_coeff\_from\_fine\_grid}).
  \item Build a symmetry-closed \texttt{bundle\_struct} with
        \texttt{R\_rot\_in\_bundle} and \texttt{sampling2bundle}
        (\texttt{gen\_bundle}).
  \item Form $\tilde V^q$ via the standard ISDF SVD/$C^{-1/2}$ path
        (\texttt{gen\_tildeVq}).
  \item Optionally validate Hartree--Fock exchange energies
        (\texttt{isdf\_validation} / \texttt{isdf\_validate\_HF}).
\end{enumerate}

\subsubsection{Code Implementation}

\begin{table}[h]
\centering
\caption{SC\_ISDF core functions (\texttt{service/+isdf/SC\_ISDF.m})}
\begin{tabular}{ll}
\toprule
Sub-function & Role \\
\midrule
\texttt{sc\_isdf\_load\_source} & Load pool id or \texttt{.mat} checkpoint \\
\texttt{sc\_isdf\_resolve\_uc\_fftgrid} & Infer $\mathbf N_{\mathrm{uc}}$ from ratio \\
\texttt{sc\_isdf\_duplicate\_sampling} & Eq.\ replication of $\mathbf R_\mu$ \\
\texttt{sc\_isdf\_rlu\_to\_lin} & Linear index on supercell FFT \\
\texttt{sc\_isdf\_build\_supercell\_object} & Assemble \texttt{isdf\_m}, set \texttt{interp\_scheme=adaptive\_sc} \\
\texttt{sc\_isdf\_build\_R\_rot\_extra} & Symmetry table on supercell \\
\bottomrule
\end{tabular}
\end{table}

Public entry point:
\begin{lstlisting}
id_sc = isdf.SC_ISDF(checkpoint_mat, [k1, k2, k3]);
\end{lstlisting}

%----------------------------------------------------------------------
\subsection{Input Configuration: \texttt{\&SUPERCELL}}
%----------------------------------------------------------------------

\subsubsection{Equations and Algorithm Description}

When \texttt{use\_sc\_isdf = .true.}, the ISDF driver skips adaptive iteration
and reads the latest \texttt{isdf\_adaptive\_checkpoint\_\{type\}\_id*.mat}
from \texttt{isdf\_source\_dir}.

\subsubsection{Code Implementation}

\begin{itemize}
  \item \texttt{input/allowed\_param\_list.m} --- field whitelist
  \item \texttt{input/default\_param\_values.m} --- defaults (\texttt{false}, $k=1$)
  \item \texttt{input/set\_default\_param\_value.m} --- validation
  \item \texttt{input/display\_input\_summary.m} --- runtime printout
\end{itemize}

Example (\texttt{Si8\_from\_Si2/test}):
\begin{lstlisting}
&SUPERCELL
  is_supercell = .true.
  use_sc_isdf = .true.
  k1 = 2
  k2 = 2
  k3 = 1
  isdf_source_dir = '../Si2/SAVE'
END &SUPERCELL
\end{lstlisting}

%----------------------------------------------------------------------
\subsection{ISDF Driver Branch}
%----------------------------------------------------------------------

\subsubsection{Equations and Algorithm Description}

Standard path: coarse grid $\to$ adaptive ISDF $\to$ \texttt{gen\_tildeVq}.

SC path:
\[
  \text{SC\_ISDF}
  \;\to\; \text{gen\_coeff\_from\_fine\_grid}
  \;\to\; \text{gen\_bundle}
  \;\to\; \text{gen\_tildeVq}
  \;\to\; \text{validate\_hf}.
\]

\subsubsection{Code Implementation}

File: \texttt{service/+isdf/driver.m}
\begin{itemize}
  \item \texttt{local\_use\_sc\_isdf(config)} --- gate on \texttt{SUPERCELL.use\_sc\_isdf}
  \item \texttt{local\_sc\_isdf\_pipeline} --- orchestrates the SC chain
  \item \texttt{local\_sc\_isdf\_checkpoint} --- locates source \texttt{.mat}
\end{itemize}

%----------------------------------------------------------------------
\subsection{Coefficient Extraction on the Supercell}
%----------------------------------------------------------------------

\subsubsection{Equations and Algorithm Description}

For each band $b$, BZ k-index, and spin:
\begin{equation}
  \texttt{coeff\_seper}(\mu,b,\mathbf k)
  = \psi_{b\mathbf k}\!\bigl(\mathbf r_{\mathrm{lin}_\mu}\bigr),
\end{equation}
where $\psi_{b\mathbf k}$ includes k-point rotation via
\texttt{wave\_functions.WF\_apply\_symm}.

\subsubsection{Code Implementation}

\texttt{service/+isdf/+coeff/gen\_coeff\_from\_fine\_grid.m}

Uses \texttt{bundle\_struct.fine\_grid\_lin} from SC\_ISDF, or falls back to
\texttt{isdf.SC\_ISDF\_r\_sampling\_to\_lin}.

%----------------------------------------------------------------------
\subsection{Symmetry Bundle Reconstruction}
%----------------------------------------------------------------------

\subsubsection{Equations and Algorithm Description}

Replicated sampling points are not automatically closed under crystal
symmetry. \texttt{gen\_bundle} collects unique orbit points on the fine FFT
grid, builds \texttt{sampling2bundle} and \texttt{R\_rot\_in\_bundle}
(bundle-index rotations), and fills \texttt{WF\_bundle}.

For \texttt{adaptive\_sc} / \texttt{coarse\_sc}, the coarse/extra split follows
the same logic as unit-cell \texttt{adaptive}.

\subsubsection{Code Implementation}

\texttt{service/+isdf/+rsymm/gen\_bundle.m}

Modifications for SC:
\begin{itemize}
  \item Recognise \texttt{adaptive\_sc}, \texttt{coarse\_sc}, \texttt{supercell} schemes
  \item Load coarse WF from \texttt{coeff\_seper} when \texttt{tmp} is empty
  \item Skip internal k-rotation self-test on supercell schemes (large \texttt{ikrot})
\end{itemize}

%----------------------------------------------------------------------
\subsection{Exchange Operator and HF Validation}
%----------------------------------------------------------------------

\subsubsection{Equations and Algorithm Description}

\texttt{gen\_tildeVq} uses \texttt{get\_u\_alpha}, which requires a populated
\texttt{bundle\_struct.R\_rot\_in\_bundle}.
HF validation compares exact HF exchange $E_x^{\mathrm{HF}}$ with ISDF
approximation $E_x^{\mathrm{ISDF}}$ over occupied/virtual band pairs in the
configured index range.

\subsubsection{Code Implementation}

\begin{itemize}
  \item \texttt{gen\_tildeVq.m} --- unchanged; called after bundle build
  \item \texttt{get\_u\_xalpha.m} --- reads \texttt{R\_rot\_in\_bundle}, \texttt{WF\_bundle}
  \item \texttt{+validation/isdf\_validation.m} --- accepts \texttt{adaptive\_sc} scheme
  \item \texttt{+validation/isdf\_validate\_HF.m} --- writes \texttt{isdf\_validate\_HF\_id*.txt}
\end{itemize}

%----------------------------------------------------------------------
\subsection{Test Layout (\texttt{test\_SC})}
%----------------------------------------------------------------------

\subsubsection{Procedure}

\begin{enumerate}
  \item Run unit-cell adaptive ISDF with \texttt{validate\_hf = .true.}
        (\texttt{energy\_band\_index\_max = 2 $\times$ occupied bands}).
  \item Save checkpoint under \texttt{./SAVE/}.
  \item For each supercell case, set \texttt{use\_sc\_isdf}, \texttt{k1/k2/k3},
        and \texttt{isdf\_source\_dir} pointing to the unit-cell \texttt{SAVE/}.
  \item Run \texttt{run\_sc\_test(case\_dir)}.
\end{enumerate}

\subsubsection{Si Test Matrix}

\begin{table}[h]
\centering
\begin{tabular}{lccc}
\toprule
Case & $N_{\mathrm{elec}}$ & occupied & \texttt{energy\_band\_index\_max} \\
\midrule
Si2 (unit)   &  8 &  4 &  8 \\
Si8 ($2\!\times\!2\!\times\!1$) & 32 & 16 & 32 \\
Si16 ($2\!\times\!2\!\times\!2$) & 64 & 32 & 64 \\
\bottomrule
\end{tabular}
\end{table}

\subsubsection{LiH Test Matrix}

Unit cell: \texttt{Materialtest/systems/LiH/1\_1\_1} ($8$ atoms, FFT $20^3$).

\begin{table}[h]
\centering
\begin{tabular}{lccc}
\toprule
Case & ratio & $N_{\mathrm{elec}}$ & \texttt{energy\_band\_index\_max} \\
\midrule
LiH\_111 (unit) & --- & 16 & 16 \\
LiH\_222\_from\_111 & $[2,2,2]$ & 128 & 128 \\
LiH\_444\_from\_111 & $[4,4,4]$ & 1024 & 1024 \\
LiH\_666\_from\_111 & $[6,6,6]$ & 3456 & 3456 \\
\bottomrule
\end{tabular}
\end{table}

%======================================================================
\section{File Index}
%======================================================================

\begin{table}[h]
\centering
\small
\begin{tabular}{ll}
\toprule
File & Status \\
\midrule
\texttt{service/+isdf/SC\_ISDF.m} & New \\
\texttt{service/+isdf/SC\_ISDF\_r\_sampling\_to\_lin.m} & New \\
\texttt{service/+isdf/+coeff/gen\_coeff\_from\_fine\_grid.m} & New \\
\texttt{service/+isdf/driver.m} & Modified \\
\texttt{service/+isdf/+rsymm/gen\_bundle.m} & Modified \\
\texttt{service/+isdf/+validation/isdf\_validation.m} & Modified \\
\texttt{input/*} (SUPERCELL block) & Modified \\
\texttt{test\_profile/test\_SC/} & New test harness \\
\bottomrule
\end{tabular}
\end{table}

\end{document}
